% ============================================================
% CAPITOLUL 8 — CONCLUZII ȘI DIRECȚII VIITOARE
% ============================================================

\chapter{Concluzii și Direcții de Dezvoltare Viitoare}
\label{ch:concluzii}

\section{Sinteza contribuțiilor}
\label{sec:sinteza_contributii}

\todosec{Rezumarea contribuțiilor principale ale lucrării:
implementările comparative, infrastructura distribuită, aplicația web.
Ce a fost atins față de obiectivele din Capitolul 1.}

\section{Concluzii privind implementările modelului SIR}
\label{sec:concluzii_sir}

\todosec{Care limbaj s-a dovedit cel mai performant și în ce condiții?
Care oferă cel mai bun compromis performanță/ușurință de utilizare?
Recomandări pentru practicieni.}

\section{Concluzii privind infrastructura distribuită}
\label{sec:concluzii_infrastructura}

\todosec{Evaluarea YugabyteDB + K8s pentru date epidemiologice:
avantaje (scalabilitate, fault tolerance), dezavantaje (complexitate operațională).
Potrivirea arhitecturii pentru sistemele e-medicine.}

\section{Limitări ale lucrării}
\label{sec:limitari}

\todosec{Limitări recunoscute: modelul SIR de bază ignoră structura de vârstă,
spațialitatea și stochasticitatea; benchmarkurile sunt limitate la o serie de
configurații hardware; interfața web nu implementează autentificare completă
în versiunea prezentată.}

\section{Direcții de dezvoltare viitoare}
\label{sec:directii_viitoare}

\todosec{Direcții:
(1) Extinderea la modele SEIR, SIRS, modele stochastice;
(2) Integrarea de modele de machine learning pentru estimarea parametrilor;
(3) Adăugarea de vizualizări spațiale (hărți de propagare);
(4) Suport pentru date în timp real (streaming cu Apache Kafka);
(5) Implementarea unui modul de alertare epidemiologică automată;
(6) Extinderea benchmarkurilor la GPU (CUDA, ROCm) pentru Julia și Python.}
